% !TEX program = pdflatex
% Example: IEEE conference paper template
% Compile with: pdflatex main.tex && bibtex main && pdflatex main.tex && pdflatex main.tex
% Requires: IEEEtran class (included in TeX Live / MiKTeX)

\documentclass[conference]{IEEEtran}

\usepackage{cite}
\usepackage{amsmath,amssymb}
\usepackage{graphicx}
\usepackage{booktabs}
\usepackage{url}

\begin{document}

\title{A Novel Approach to Signal Processing
       Using Deep Learning}

\author{
\IEEEauthorblockN{Author One}
\IEEEauthorblockA{Department, University\\
City, Country\\
Email: one@uni.edu}
\and
\IEEEauthorblockN{Author Two}
\IEEEauthorblockA{Department, University\\
City, Country\\
Email: two@uni.edu}
}

\maketitle

\begin{abstract}
This paper proposes a novel approach to signal processing
using deep learning. We introduce a multi-scale convolutional
architecture that captures both local and global features.
\end{abstract}

\begin{IEEEkeywords}
signal processing, deep learning, neural networks.
\end{IEEEkeywords}

\section{Introduction}
\IEEEPARstart{S}{ignal} processing has traditionally relied
on hand-crafted features. Deep learning offers an alternative
that learns features directly from data.

\section{Related Work}
Prior work~\cite{ref1,ref2} has explored CNN-based signal
processing, but limited to single-scale architectures.

\section{Proposed Method}
Our multi-scale architecture processes the input signal
at three resolutions simultaneously:
\begin{equation}
y_i = \text{Conv}_{k_i}(x), \quad i \in \{3, 5, 7\}
\end{equation}

\section{Experiments}
Table~\ref{tab:results} shows the performance comparison.

\begin{table}[htbp]
\caption{Performance comparison.}
\label{tab:results}
\centering
\begin{tabular}{lcc}
\toprule
Method & Accuracy & Latency \\
\midrule
Baseline & 85.2\% & 12ms \\
Single-scale & 89.5\% & 18ms \\
Proposed & 93.7\% & 22ms \\
\bottomrule
\end{tabular}
\end{table}

\section{Conclusion}
Our multi-scale approach achieves 93.7\% accuracy with
acceptable latency for real-time applications.

\bibliographystyle{IEEEtran}
\bibliography{references}

\end{document}
